\documentclass[10pt,a4paper]{amsart}
\usepackage{amsmath,amssymb,amsthm,graphicx}
\usepackage{verbatim}
\newenvironment{reference}{\comment}{\endcomment}
%\newenvironment{reference}{}{}
\newenvironment{slogan}{\comment}{\endcomment}
\newenvironment{history}{\comment}{\endcomment}

% For commutative diagrams we use Xy-pic
\usepackage[all]{xy}

% We use 2cell for 2-commutative diagrams.
\xyoption{2cell}
% Theorem environments.
%
\theoremstyle{plain}
\newtheorem{theorem}[subsection]{Theorem}
\newtheorem{proposition}[subsection]{Proposition}
\newtheorem{lemma}[subsection]{Lemma}

\theoremstyle{definition}
\newtheorem{definition}[subsection]{Definition}
\newtheorem{example}[subsection]{Example}
\newtheorem{exercise}[subsection]{Exercise}
\newtheorem{situation}[subsection]{Situation}

\theoremstyle{remark}
\newtheorem{remark}[subsection]{Remark}
\newtheorem{remarks}[subsection]{Remarks}

\numberwithin{equation}{subsection}

% Macros
%
\def\lim{\mathop{\mathrm{lim}}\nolimits}
\def\colim{\mathop{\mathrm{colim}}\nolimits}
\def\Spec{\mathop{\mathrm{Spec}}}
\def\Hom{\mathop{\mathrm{Hom}}\nolimits}
\def\Ext{\mathop{\mathrm{Ext}}\nolimits}
\def\SheafHom{\mathop{\mathcal{H}\!\mathit{om}}\nolimits}
\def\SheafExt{\mathop{\mathcal{E}\!\mathit{xt}}\nolimits}
\def\Sch{\mathit{Sch}}
\def\Mor{\mathop{\mathrm{Mor}}\nolimits}
\def\Ob{\mathop{\mathrm{Ob}}\nolimits}
\def\Sh{\mathop{\mathit{Sh}}\nolimits}
\def\NL{\mathop{N\!L}\nolimits}
\def\CH{\mathop{\mathrm{CH}}\nolimits}
\def\proetale{{pro\text{-}\acute{e}tale}}
\def\etale{{\acute{e}tale}}
\def\QCoh{\mathit{QCoh}}
\def\Ker{\mathop{\mathrm{Ker}}}
\def\Im{\mathop{\mathrm{Im}}}
\def\Coker{\mathop{\mathrm{Coker}}}
\def\Coim{\mathop{\mathrm{Coim}}}

% Boxtimes
%
\DeclareMathSymbol{\boxtimes}{\mathbin}{AMSa}{"02}

%
% Macros for moduli stacks/spaces
%
\def\QCohstack{\mathcal{QC}\!\mathit{oh}}
\def\Cohstack{\mathcal{C}\!\mathit{oh}}
\def\Spacesstack{\mathcal{S}\!\mathit{paces}}
\def\Quotfunctor{\mathrm{Quot}}
\def\Hilbfunctor{\mathrm{Hilb}}
\def\Curvesstack{\mathcal{C}\!\mathit{urves}}
\def\Polarizedstack{\mathcal{P}\!\mathit{olarized}}
\def\Complexesstack{\mathcal{C}\!\mathit{omplexes}}
% \Pic is the operator that assigns to X its picard group, usage \Pic(X)
% \Picardstack_{X/B} denotes the Picard stack of X over B
% \Picardfunctor_{X/B} denotes the Picard functor of X over B
\def\Pic{\mathop{\mathrm{Pic}}\nolimits}
\def\Picardstack{\mathcal{P}\!\mathit{ic}}
\def\Picardfunctor{\mathrm{Pic}}
\def\Deformationcategory{\mathcal{D}\!\mathit{ef}}
\usepackage{cleveref}
\crefname{originalreference}{original reference}{original references}
\Crefname{originalreference}{Original reference}{Original references}
\makeatletter
\newlabel{lemma-intersect-normal}{{09HP}{source}{}{sourcepointer-1}{}}
\newlabel{lemma-intersect-normal@cref}{{[originalreference][0][]09HP}{[1][0][]source}{}{}{}}
\makeatother
\crefname{originalreference}{original reference}{original references}
\Crefname{originalreference}{Original reference}{Original references}
\makeatletter
\newlabel{lemma-map-into-algebraic-closure}{{09GU}{source}{}{sourcepointer-2}{}}
\newlabel{lemma-map-into-algebraic-closure@cref}{{[originalreference][0][]09GU}{[1][0][]source}{}{}{}}
\makeatother
\title{Existence and uniqueness of the normal closure of a finite extension}
\author{The Stacks Project Authors}
\date{stacks-project:fields.tex:local-snapshot-20261012; GNU FDL 1.2 or later}
\begin{document}
\maketitle
\noindent\textit{Editorial note.} All delivered mathematical text is the author's own, copied from the pinned Stacks Project chapter file fields.tex; no statement or proof step was rewritten, reordered or removed.\par\medskip
\noindent\textit{Editorial note.} The counted proof is the author's complete argument: The following definition-normal-closure (source lines 2340-2345) names the field constructed here; the inside-a-given-normal-extension variant is lemma-normal-closure-inside-normal later in the chapter.\par\medskip
\noindent\textit{Editorial note.} The label lemma-normal-closure names the counted statement, and the proof environment immediately following it in the source is delivered verbatim as its proof.\par\medskip
\noindent\textit{Editorial note.} The counted statement is carried by the author's own lemma environment in the chapter source.\par\medskip
\noindent\textit{Allowed standard prerequisites (editorial).}\begin{itemize}
\item Same chapter: lemma-splitting-field (smallest splitting extension exists, is normal and unique up to isomorphism, source lines 2265-2293), lemma-map-into-algebraic-closure (embeddings into an algebraic closure extend along algebraic extensions), lemma-intersect-normal (the intersection of two normal subextensions in a common overfield is normal), lemma-algebraically-closed (factorization over an algebraic closure)
\end{itemize}\medskip
\ifdefined\hypertarget\hypertarget{sourcepointer-1}{}\fi
\noindent\textit{Source reference.} Original reference 09HP: sources/stacks-fields/pinned-original-source.tex, source lines 2068--2068. Prerequisite pointer invoked at 2333; it names a standard fact of the same chapter that the counted argument uses, and every step of the counted proof built on it is retained verbatim in the delivered body.\par\medskip
\ifdefined\hypertarget\hypertarget{sourcepointer-2}{}\fi
\noindent\textit{Source reference.} Original reference 09GU: sources/stacks-fields/pinned-original-source.tex, source lines 1240--1240. Prerequisite pointer invoked at 2330; it names a standard fact of the same chapter that the counted argument uses, and every step of the counted proof built on it is retained verbatim in the delivered body.\par\medskip
\begin{lemma}
\label{lemma-normal-closure}
\begin{slogan}
Existence of normal closure of finite extensions of fields.
\end{slogan}
Let $E/F$ be a finite extension of fields. There exists a unique
smallest finite extension $K/E$ such that $K$ is normal over $F$.
\end{lemma}

\begin{proof}
Choose generators $\alpha_1, \ldots, \alpha_n$ of $E$ over $F$.
Let $P_1, \ldots, P_n$ be the minimal polynomials of
$\alpha_1, \ldots, \alpha_n$ over $F$. Set $P = P_1 \ldots P_n$.
Observe that $(x - \alpha_1) \ldots (x - \alpha_n)$ divides $P$, since
each $(x - \alpha_i)$ divides $P_i$. Say
$P = (x - \alpha_1) \ldots (x - \alpha_n)Q$.
Let $K/E$ be the splitting field of $P$ over $E$.
We claim that $K$ is the splitting field of $P$ over $F$ as well
(which implies that $K$ is normal over $F$).
This is clear because $K/E$ is generated by the roots of
$Q$ over $E$ and $E$ is generated by the roots of
$(x - \alpha_1) \ldots (x - \alpha_n)$ over $F$, hence
$K$ is generated by the roots of $P$ over $F$.

\medskip\noindent
Uniqueness. Suppose that $K'/E$ is a second smallest extension such that
$K'/F$ is normal. Choose an algebraic closure $\overline{F}$ and an
embedding $\sigma_0 : E \to \overline{F}$. By
Lemma \ref{lemma-map-into-algebraic-closure}
we can extend $\sigma_0$ to $\sigma : K \to \overline{F}$ and
$\sigma' : K' \to \overline{F}$.
By Lemma \ref{lemma-intersect-normal} we see that
$\sigma(K) \cap \sigma'(K')$ is normal over $F$.
By minimality we conclude that $\sigma(K) = \sigma'(K')$.
Thus $\sigma^{-1} \circ \sigma' : K' \to K$ gives an isomorphism
of extensions of $E$.
\end{proof}

\end{document}
